\section{Introduction}

Kepler, building on Tycho's observations, first extracted empirical regularities: planets move on ellipses with the Sun at a focus, sweep out equal areas in equal times, and across planets satisfy the invariant period--size relation $T^2\propto a^3$ \cite{ostermannwanner2012geometry}.

In Book~I of the \emph{Principia} (1687), Newton formulates the key hinge in two parts. Proposition~I (Theorem~I) and Proposition~II (Theorem~II) establish the equivalence between central direction and area sweep proportional to time in planar motion; with the conic-focus condition added, Newton then derives inverse-square dependence \cite{newton1846mottewikisource,chandrasekhar1995principia}. In the first-edition tradition, the conic-specific historical forward chain is commonly cited as Propositions~XI--XIII, while the complementary historical reverse reconstruction chain is often cited through Propositions~XVII and~XLI. Unless noted otherwise, proposition/theorem references in this paper follow the 1726 third-edition sequence as presented in the Motte--Cajori translation line \cite{newton1846mottewikisource,newton1934mottecajori}. Chandrasekhar's common-reader commentary is widely associated with this translation tradition \cite{chandrasekhar1995principia,newton1934mottecajori}.

Historically, orbit-to-force was the direction of discovery, while force-to-orbit was its complement. In the terminology used in this paper, these are the inverse problem and forward problem, respectively. Newton's geometric arguments in Book~I treat both directions for ellipse, parabola, and hyperbola through conic-specific propositions. Under inverse-square centripetal attraction, trajectories are conic sections with the force center at a focus, and the orbit type (ellipse/parabola/hyperbola) is determined by the motion regime (equivalently, by the constants fixed by initial data). In revisions prepared around 1712 and published in the second edition (1713), Newton also included additional forward-facing comparisons that refer motion to the conic center (center of ellipse or hyperbola), complementing the focus-centered constructional route \cite{Nov95,chandrasekhar1995principia}. Modern reconstructions of Newton's argument structure are given in \cite{chandrasekhar1995principia,markowsky2011newton}.

The continental differential-calculus tradition associated with Leibniz and Bernoulli emphasized analytic generality, while Newton's treatment remained geometric and limit-based (first and last ratios). Johann Bernoulli and others criticized parts of Newton's forward-problem argument, and this methodological dispute shaped later celestial mechanics \cite{guicciardini1999reading}. A modern defense of Newton's internal deductive structure is given by Chandrasekhar (especially pp.~111--112), and a related historical assessment is given by Arnol'd \cite{chandrasekhar1995principia,arnold1990huygens}. In modern terms, fully explicit ODE existence/uniqueness theorems came later, especially with 19th-century work led by Cauchy.

Modern pedagogy usually starts from Newton's laws and a force model as given, then derives orbital laws as consequences. In the rest of this paper, we use forward/inverse in this modern sense unless explicitly marked as historical. Thus Kepler-to-$1/r^2$ is the modern inverse direction, while $1/r^2$-to-conic is the modern forward direction. Standard modern routes for the former are summarized in \cite{chandrasekhar1995principia,ostermannwanner2012geometry,markowsky2011newton}; for the latter, a common pedagogical line runs through Hamilton (1847), Maxwell (1877), and Feynman's 1964 ``lost lecture,'' and is further refined in vHH (2009) and CRS (2016) \cite{hamilton1847hodograph,maxwell1877mattermotion,goodstein1996feynman,vanhaandelheckman2009teaching,carinena2016newlook}. Our aim is to present geometric proofs in both directions, in a style close to Newton's Euclidean constructive approach.

We briefly summarize what the main cited references emphasize, to clarify how the present proof fits among them (we refer to the bibliography for full details). For historical discussion of how Book~I evolves across different versions/editions, see also \cite{Nov95,chandrasekhar1995principia}.
\begin{itemize}

  \item Nobel laureate Chandrasekhar (1995) \cite{chandrasekhar1995principia}. A detailed modern reconstruction of Newton's arguments, combining synthetic geometry with contemporary analytic language. Some derivations do not follow Newton's original route step by step, so this book is best read together with the \emph{Principia} itself. Nonetheless, it is indispensable and greatly aids reading the \emph{Principia}.

  \item Guicciardini (1999) \cite{guicciardini1999reading}. A historical study of the post-\emph{Principia} debate on Newton's mathematical method, including Bernoulli-era criticism and the geometry-vs-analysis fault line.

  \item Arnol'd (1990) \cite{arnold1990huygens}. A historical-mathematical perspective on the Newton--Hooke/Huygens--Barrow line, often cited in discussions of how Newton's geometric method should be interpreted against later analytic standards.

  \item Sir William Rowan Hamilton (1847) \cite{hamilton1847hodograph}. Hamilton introduces the hodograph and proves circular hodographs for inverse-square central attraction. This is a landmark exposition, notable for its distinctive communication style: almost no diagrams and very light symbolic machinery, yet a strikingly clear geometric argument. Hamilton is an early, prominent British-Isles figure in the 19th-century analytical reform and re-connection with continental methods, while still keeping a strongly synthetic/geometric mode of thought; this work is a prime example of that style.

  \item Maxwell (1877) \cite{maxwell1877mattermotion}. In this introductory treatment of motion and matter, Maxwell develops the pedagogical hodograph viewpoint and emphasizes its value for reconstructing orbital geometry from inverse-square dynamics. He gives an early clear construction in which a rotated hodograph circle corresponds to the ellipse's directrix circle, making the conic-to-inverse-square (modern inverse problem) derivation especially transparent for students.

  \item Goodstein \& Goodstein (1996) \cite{goodstein1996feynman}. A geometric narrative inspired by Feynman's 1964 ``lost lecture,'' emphasizing a nearly calculus-free route for the forward problem (inverse-square $\Rightarrow$ conic orbits) using hodograph synthesis.

  \item Derbes (2001) \cite{derbes2001reinventing}. A clear pedagogical account of hodographic (velocity-space) methods for the Kepler problem, including the parabolic case.

  \item Stavek (2019) \cite{stavek2019newtonsparabola}. A geometric exploration centered on Newton's parabola and related classical constructions (directrix, pedal curve, evolute, subtangent/subnormal, and a Ptolemy-circle/hodograph viewpoint).

  \item Markowsky (2011) \cite{markowsky2011newton}. A careful retelling and streamlining of Newton's arguments, with an emphasis on geometric structure and pedagogy. Markowsky also emphasizes that whether Newton regarded ``having a solution'' as sufficient for the forward problem (inverse-square $\Rightarrow$ conics) is debatable: one can justify the step via existence/uniqueness results for the associated ordinary differential equations. It also seeks a \emph{Principia}-style proof of energy conservation, from which one can derive that the trajectory is a conic.

  \item Provost \& Bracco (2008) \cite{provostbracco2008simple}. A hybrid yet insightful and concise derivation highlighting conserved quantities---essentially the Laplace--Runge--Lenz vector---that analytical approaches often use to show that the path is a conic. This paper offers an elementary view based on areal speed and related geometric relations, and also summarizes Jakob Hermann's early-18th-century proof (c.\ 1710s) and the aforementioned proof of Hamilton (1847).

  \item Woan \cite{a1dynamicalastronomy}. A standard mechanics approach in polar coordinates, deriving the conic form by solving for $r(\theta)$ using angular momentum and energy, in the tradition of analytic approaches developed by mathematicians such as Bernoulli (late 17th to early 18th century).

  \item Simha (2021) \cite{simha2021algebra}. A calculus-light, trigonometric presentation of Kepler's first law; it substantially overlaps with the Markowsky-style approach, possibly without awareness of that earlier source. It is mentioned here because it uses the law of cosines effectively.

  \item van~Haandel \& Heckman (2009) \cite{vanhaandelheckman2009teaching}. A freshman-oriented teaching note emphasizing geometric intuition. Besides giving a quick account of Feynman and Newton's methods, it also provides a geometric construction of the LRL vector that offers further insight into the forward problem.

  \item Cari\~nena--Ra\~nada--Santander (2016) \cite{carinena2016newlook}. A modern refinement of the hodograph/Feynman line (referred to in this paper as CRS16), including a cleaner force-center transformation framework.
\end{itemize}

In Section~2 we briefly repeat Newton's argument (with his classic figure) showing why the area theorem implies a centripetal (central) force. In Sections~3--5 we present geometric proofs for inverse-square dependence using the auxiliary circle, affine transport, and conic-specific refinements. Section~6 then treats the complementary forward problem through a translated, $\Delta t$-scaled hodograph construction.
